\pagestyle{empty}
\tight
\begin{tblr}{width=\textwidth,colspec={|Q[c,m,1.9cm]|X[l,m]|},hlines,vlines,hspan=minimal,row{1}={bg=headblue},rowsep=1pt,colsep=3pt}
\SetCell{c,m}\hfil\logo\hfil & \textbf{POLITEKNIK NEGERI MALANG}\par\textbf{JURUSAN TEKNOLOGI INFORMASI}\par\textbf{PROGRAM STUDI: D4 TEKNIK INFORMATIKA} \\
\SetCell[c=2]{c} \textbf{RENCANA PEMBELAJARAN SEMESTER (RPS) --- DRAF KONSEP A: SOLUSI TI} & \\
\end{tblr}
\begin{longtblr}{width=\textwidth,colspec={|Q[l,m,1.9cm]|X[0.85,l,m]|X[1.45,l,m]|X[0.75,l,m]|X[1.15,l,m]|},hlines,vlines,hspan=minimal,rowsep=1pt,colsep=2pt,row{1,3}={bg=headgray,font=\bfseries}}
MATA KULIAH & KODE & BOBOT (SKS)/jam & SEMESTER & TGL. PENYUSUNAN \\
Kuliah Kerja Nyata Tematik & RTI257002 (sementara)\footnotemark[1] & 2 SKS / 4 jam/minggu & 6/7 (sementara)\footnotemark[1] & 13 Juli 2026 \\
OTORISASI & Dosen Pengembang RPS & Koordinator RMK & \SetCell[c=2]{l} Ka PRODI &  \\
 & Tim Pengajar Program Studi D4 TI &  & \SetCell[c=2]{l}{\vspace{8mm}\par Dr. Ely Setyo Astuti, S.T., M.T.} &  \\
\SetCell[r=14]{l} Capaian Pembelajaran (CP) & \SetCell[c=4]{l,bg=headgray,font=\bfseries} Capaian Pembelajaran Lulusan Program Studi (CPL-Prodi) &  &  &  \\
 & CPL01 & \SetCell[c=3]{l} Menginternalisasi ketakwaan kepada Tuhan Yang Maha Esa, menaati hukum, disiplin, dan menunjukkan profesionalisme melalui etika profesi, pembelajaran sepanjang hayat, serta respons terhadap isu sosial dan teknologi. &  &  \\
 & CPL03 & \SetCell[c=3]{l} Mampu mengelola tim, manajemen diri, berkomunikasi secara lisan maupun tertulis dengan baik dan mampu melakukan presentasi. &  &  \\
 & CPL06 & \SetCell[c=3]{l} Mampu merancang, mengimplementasikan, menguji, melakukan deployment, dan memelihara, serta menjamin mutu perangkat lunak yang menjawab permasalahan dan memenuhi kebutuhan stakeholder. &  &  \\
 & CPL08 & \SetCell[c=3]{l} Mampu mengelola proyek teknologi informasi secara profesional dengan mengoptimalkan sumber daya yang tersedia. &  &  \\
 & \SetCell[c=4]{l,bg=headgray,font=\bfseries} Capaian Pembelajaran mata kuliah (CPL-MK) &  &  &  \\
 & CPMK0101 & \SetCell[c=3]{l} Mampu menunjukkan profesionalisme melalui penerapan etika profesi TI, keselamatan kerja, kewirausahaan bertanggung jawab, rekayasa perangkat lunak ramah lingkungan, dan disiplin dalam penelitian maupun penerapan aplikasi teknologi. &  &  \\
 & CPMK0301 & \SetCell[c=3]{l} Mampu mengelola diri secara profesional, berkomunikasi secara lisan maupun tertulis, dan melakukan presentasi secara efektif dalam konteks kerja. &  &  \\
 & CPMK0803 & \SetCell[c=3]{l} Mampu melaksanakan, memonitor, dan mengendalikan proyek serta mengelola dinamika tim secara efektif. &  &  \\
 & CPMK0607 & \SetCell[c=3]{l} Mampu membangun solusi sistem informasi terintegrasi yang kompleks dan menjawab kebutuhan \textit{stakeholder}. &  &  \\
 & \SetCell[c=4]{l,bg=headgray,font=\bfseries} Kemampuan akhir tiap tahapan belajar (Sub-CPMK) &  &  &  \\
 & SCPMK0101-06001 & \SetCell[c=3]{l} Mahasiswa mampu menunjukkan etika kerja dan profesionalisme selama observasi dan pelaksanaan program kerja di lokasi mitra. &  &  \\
 & SCPMK0301-06002 & \SetCell[c=3]{l} Mahasiswa mampu menggali kebutuhan mitra secara partisipatif dan mengomunikasikan hasil program kerja dan solusi TI. &  &  \\
 & SCPMK0803-06003 & \SetCell[c=3]{l} Mahasiswa mampu menyusun dan melaksanakan rencana program kerja KKN sesuai jadwal yang disepakati bersama mitra. &  &  \\
 & SCPMK0607-06004 & \SetCell[c=3]{l} Mahasiswa mampu merancang, mengimplementasikan, dan menyerahterimakan solusi TI yang sesuai kebutuhan mitra. &  &  \\
\SetCell[r=2]{l} Penilaian & \SetCell[c=4]{l} *Nilai CPMK didapat dari agregasi nilai SUB-CPMK yang dibebankan &  &  &  \\
 & \SetCell[c=4]{l}{\tiny\begin{tblr}{width=\linewidth,colspec={|Q[c,m,0.1\linewidth]|Q[c,m,0.16\linewidth]|X[l,m]|X[l,m]|X[l,m]|X[l,m]|Q[c,m,0.08\linewidth]|},hlines,vlines,hspan=minimal,rowsep=1pt,colsep=0.7pt,row{1,2}={bg=headgray,font=\bfseries}}
Kode CPMK & Kode SCPMK & \SetCell[c=4]{c} Bobot per Bentuk Penilaian & & & & Total Bobot \\
 & & Proposal & Produk/Solusi TI & Laporan & Seminar & \\
CPMK0101 & SCPMK0101-06001 & 5\%: etika \& kelayakan program kerja & & 5\%: refleksi etika \& tanggung jawab profesional & & 10\% \\
CPMK0301 & SCPMK0301-06002 & 5\%: komunikasi kebutuhan dengan mitra & & & 10\%: komunikasi hasil kepada mitra dan penguji & 15\% \\
CPMK0803 & SCPMK0803-06003 & 5\%: perencanaan \& penjadwalan program kerja & 15\%: pelaksanaan \& pengelolaan program & 10\%: dokumentasi pelaksanaan \& evaluasi & & 30\% \\
CPMK0607 & SCPMK0607-06004 & & 25\%: rancang bangun \& implementasi solusi TI & 10\%: dokumentasi teknis solusi & 10\%: defense solusi teknis & 45\% \\
\SetCell[c=6]{r}\textbf{Total Per Penilaian} & & 15\% & 40\% & 25\% & 20\% & \textbf{100\%} \\
\end{tblr}} &  &  &  \\
Diskripsi Singkat Mata Kuliah & \SetCell[c=4]{l} Kuliah Kerja Nyata Tematik (KKN Tematik) adalah mata kuliah wajib 2 SKS penunjang IKU kegiatan mahasiswa di luar Prodi (dasar: Nota Dinas Wakil Direktur I No. 11/WADIR I/KM/2026 dan Permendiktisaintek No. 39 Tahun 2025) yang menempatkan mahasiswa pada mitra/komunitas untuk merancang dan mengimplementasikan solusi teknologi informasi yang menjawab kebutuhan nyata mitra, disertai laporan dan seminar hasil. \\
Bahan Kajian / Materi Pembelajaran & \SetCell[c=4]{l}\begin{enumerate}\item Pembekalan KKN, etika kerja lapangan, dan pemetaan mitra/komunitas\item Analisis kebutuhan mitra dan perancangan solusi TI\item Pelaksanaan program kerja dan implementasi solusi TI di lokasi mitra\item Pelatihan pengguna dan evaluasi dampak solusi TI\item Penyusunan laporan KKN dan seminar hasil\end{enumerate} \\
\SetCell[r=2]{l} Pustaka & \SetCell[c=4]{l} \textbf{Utama:}\begin{enumerate}\item Buku Pedoman KKN Tematik Politeknik Negeri Malang.\item Panduan MBKM Kemdikbud 2020.\item Pedoman Penulisan Laporan KKN Polinema.\end{enumerate} \\
 & \SetCell[c=4]{l} \textbf{Pendukung:} Materi LMS, dokumen profil mitra/komunitas, dan bahan bimbingan dosen pembimbing lapangan (DPL). \\
Media Pembelajaran & \SetCell[c=2]{l} \textbf{Software:} Aplikasi logbook digital, LMS, perangkat bantu pengembangan solusi TI. &  & \SetCell[c=2]{l} \textbf{Hardware:} Komputer/laptop, perangkat lapangan mitra. &  \\
Nama Dosen Pengampu & \SetCell[c=4]{l} Tim Pengajar Program Studi D4 TI \\
Matakuliah Syarat & \SetCell[c=4]{l} Lulus minimal 90 SKS \\
\end{longtblr}
\begin{longtblr}{width=\textwidth,colspec={|Q[c,m,0.06\textwidth]|X[1.35,l,m]|X[1.08,l,m]|X[1.16,l,m]|X[0.7,l,m]|X[1.24,l,m]|X[1.06,l,m]|X[0.98,l,m]|Q[c,m,0.05\textwidth]|},hlines,vlines,hspan=minimal,rowhead=2,row{1,2}={bg=headgreen,font=\bfseries},rowsep=1pt,colsep=1.5pt}
Mgg.\par Ke & Kemampuan Akhir Yang Direncanakan (Sub-CP-MK) & Bahan kajian (Materi Pembelajaran) & Bentuk dan Metode Pembelajaran & Estimasi Waktu & Pengalaman Belajar Mahasiswa & Kriteria \& Bentuk Penilaian & Indikator Penilaian & Bobot Penilaian (\%) \\
(1) & (2) & (3) & (4) & (5) & (6) & (7) & (8) & (9) \\
1 & SCPMK0101-06001, SCPMK0803-06003 & Pembekalan KKN: orientasi, etika kerja lapangan, dan kontrak belajar. & Modalitas: Blended Learning. Bentuk: Luring/Daring. Metode: Case Method, ceramah interaktif, dan diskusi. & 1 x 2 x 50' tatap muka; persiapan mandiri. & Mahasiswa mengikuti pembekalan KKN, memahami etika kerja di lokasi mitra, dan menandatangani kontrak belajar. & Bentuk penilaian: Kehadiran dan partisipasi pembekalan. Mengacu rubrik RTM. & Ketepatan pemahaman prosedur; kesiapan etika kerja; kelengkapan kontrak belajar. & 2 \\
2 & SCPMK0301-06002, SCPMK0803-06003 & Observasi lapangan: pemetaan mitra/komunitas dan potensi wilayah. & Modalitas: Praktik Lapangan dan Bimbingan. Metode: observasi lapangan, wawancara mitra, dan bimbingan dosen pembimbing lapangan (DPL). & 1 x 2 x 50' observasi lapangan; kerja mandiri. & Mahasiswa melakukan observasi ke lokasi mitra/komunitas dan mengidentifikasi potensi masalah yang dapat diselesaikan dengan solusi TI. & Bentuk penilaian: Catatan observasi lapangan. Mengacu rubrik RTM. & Ketepatan identifikasi masalah; kelengkapan data lapangan. & 2 \\
3 & SCPMK0301-06002, SCPMK0607-06004 & Analisis kebutuhan mitra dan spesifikasi awal solusi TI. & Modalitas: Praktik Lapangan dan Bimbingan. Metode: analisis partisipatif dan bimbingan DPL. & 1 x 2 x 50' bimbingan; analisis mandiri. & Mahasiswa menganalisis kebutuhan mitra secara partisipatif dan merumuskan spesifikasi awal solusi TI yang akan dibangun. & Bentuk penilaian: Draft analisis kebutuhan. Mengacu rubrik RTM. & Ketepatan analisis kebutuhan; kesesuaian solusi dengan konteks mitra. & 3 \\
4 & SCPMK0101-06001, SCPMK0803-06003 & Penyusunan proposal program kerja KKN, termasuk rancangan solusi TI. & Modalitas: Blended Learning. Bentuk: Luring/Daring. Metode: Case Method, penyusunan dokumen, dan bimbingan DPL. & 1 x 2 x 50' bimbingan; penyusunan mandiri. & Mahasiswa menyusun proposal program kerja KKN yang memuat rencana kegiatan, jadwal, dan rancangan solusi TI bagi mitra. & Bentuk penilaian: Draft proposal program kerja. Mengacu rubrik RTM. & Kelengkapan proposal; kelayakan rencana kerja; ketepatan etika penyusunan. & 5 \\
5 & SCPMK0101-06001, SCPMK0301-06002, SCPMK0803-06003 & Seminar proposal program kerja KKN. & Modalitas: Blended Learning. Bentuk: Luring/Daring. Metode: presentasi dan Case Method. & 1 x 2 x 50' seminar proposal. & Mahasiswa mempresentasikan proposal program kerja KKN di hadapan DPL dan menerima masukan untuk perbaikan. & Bentuk penilaian: Seminar proposal program kerja. Mengacu rubrik RTM. & Kelayakan program kerja; kejelasan rancangan solusi TI; kualitas presentasi. & 15 \\
6 & SCPMK0803-06003, SCPMK0607-06004 & Pelaksanaan program kerja tahap 1: pengumpulan data dan desain solusi TI. & Modalitas: Praktik Lapangan dan Bimbingan. Metode: praktik lapangan, monitoring berkala, dan bimbingan DPL. & 1 x 2 x 50' bimbingan; kerja lapangan mandiri. & Mahasiswa melaksanakan tahap awal program kerja di lokasi mitra dan mendesain solusi TI sesuai kebutuhan yang telah dianalisis. & Bentuk penilaian: Logbook dan draft desain solusi. Mengacu rubrik RTM. & Ketepatan waktu; kesesuaian desain; kualitas dokumentasi. & 3 \\
7 & SCPMK0607-06004 & Pelaksanaan program kerja tahap 1 (lanjutan): implementasi awal solusi TI. & Modalitas: Praktik Lapangan dan Bimbingan. Metode: praktik lapangan, monitoring berkala, dan bimbingan DPL. & 1 x 2 x 50' bimbingan; kerja lapangan mandiri. & Mahasiswa mengimplementasikan bagian awal solusi TI dan menguji cobakannya bersama mitra. & Bentuk penilaian: Logbook dan progress implementasi. Mengacu rubrik RTM. & Kualitas implementasi; ketepatan waktu; keterlibatan mitra. & 3 \\
8 & SCPMK0101-06001, SCPMK0803-06003 & Monitoring dan evaluasi (monev) tengah bersama mitra dan DPL. & Modalitas: Praktik Lapangan dan Bimbingan. Metode: praktik lapangan, monitoring berkala, dan bimbingan DPL. & Disesuaikan jadwal mitra. & Mahasiswa menjalani monev tengah program KKN bersama mitra dan DPL untuk mengevaluasi kemajuan program kerja dan solusi TI. & Bentuk penilaian: Formulir monev tengah oleh mitra/DPL. Mengacu rubrik RTM. & Kesesuaian progres dengan rencana; profesionalisme; kualitas kemajuan solusi. & 15 \\
9 & SCPMK0607-06004 & Pelaksanaan program kerja tahap 2: penyempurnaan solusi TI berdasarkan masukan monev. & Modalitas: Praktik Lapangan dan Bimbingan. Metode: praktik lapangan, monitoring berkala, dan bimbingan DPL. & 1 x 1 x 50' monitoring; kerja mandiri. & Mahasiswa menyempurnakan solusi TI berdasarkan masukan hasil monev tengah. & Bentuk penilaian: Logbook dan revisi solusi. Mengacu rubrik RTM. & Kualitas revisi; ketepatan waktu. & 3 \\
10 & SCPMK0301-06002, SCPMK0607-06004 & Pelaksanaan program kerja tahap 2 (lanjutan): pelatihan penggunaan solusi TI kepada mitra. & Modalitas: Praktik Lapangan dan Bimbingan. Metode: praktik lapangan, monitoring berkala, dan bimbingan DPL. & 1 x 1 x 50' monitoring; kerja mandiri. & Mahasiswa melatih mitra menggunakan solusi TI yang dibangun dan mengumpulkan umpan balik pengguna. & Bentuk penilaian: Logbook dan berita acara pelatihan. Mengacu rubrik RTM. & Kualitas pelatihan; ketepatan waktu; keterlibatan mitra. & 3 \\
11 & SCPMK0803-06003 & Pelaksanaan program kerja tahap 3: evaluasi dampak program kepada komunitas. & Modalitas: Praktik Lapangan dan Bimbingan. Metode: praktik lapangan, monitoring berkala, dan bimbingan DPL. & 1 x 1 x 50' monitoring; kerja mandiri. & Mahasiswa mengevaluasi dampak program kerja dan solusi TI terhadap komunitas/mitra. & Bentuk penilaian: Logbook dan catatan evaluasi dampak. Mengacu rubrik RTM. & Ketepatan waktu; kualitas evaluasi dampak. & 3 \\
12 & SCPMK0607-06004 & Konsolidasi luaran dan dokumentasi teknis solusi TI. & Modalitas: Praktik Lapangan dan Bimbingan. Metode: praktik lapangan, monitoring berkala, dan bimbingan DPL. & 1 x 1 x 50' monitoring; kerja mandiri. & Mahasiswa menyusun dokumentasi teknis solusi TI (manual pengguna, dokumentasi sistem) sebagai luaran program. & Bentuk penilaian: Dokumentasi teknis solusi TI. Mengacu rubrik RTM. & Kelengkapan dokumentasi; ketepatan waktu. & 3 \\
13 & SCPMK0101-06001, SCPMK0301-06002 & Bimbingan penyusunan laporan KKN. & Modalitas: Blended Learning. Bentuk: Luring/Daring. Metode: Case Method dan bimbingan DPL. & 1 x 1 x 50' bimbingan; penulisan mandiri. & Mahasiswa mendapatkan bimbingan teknis penulisan laporan KKN mencakup sistematika, konten, dan tata tulis ilmiah. & Bentuk penilaian: Draft laporan dengan catatan bimbingan. Mengacu rubrik RTM. & Ketepatan waktu; kualitas dokumen. & 5 \\
14 & SCPMK0301-06002 & Finalisasi laporan KKN. & Modalitas: Blended Learning. Bentuk: Luring/Daring. Metode: Case Method dan bimbingan DPL. & 1 x 1 x 50' bimbingan; penulisan mandiri. & Mahasiswa menyelesaikan draft akhir laporan KKN mencakup seluruh bab dan lampiran. & Bentuk penilaian: Laporan KKN final. Mengacu rubrik RTM. & Kelengkapan laporan; ketepatan waktu. & 5 \\
15 & SCPMK0101-06001, SCPMK0803-06003 & Persiapan seminar hasil: penyusunan bahan presentasi dan kelengkapan dokumen. & Modalitas: Blended Learning. Bentuk: Luring/Daring. Metode: Case Method dan bimbingan DPL. & 1 x 1 x 50' persiapan; mandiri. & Mahasiswa menyiapkan bahan presentasi seminar hasil dan melengkapi dokumen pendukung program kerja. & Bentuk penilaian: Kelengkapan dokumen seminar. Mengacu rubrik RTM. & Kelengkapan dokumen; kesiapan presentasi. & 5 \\
16 & SCPMK0301-06002, SCPMK0607-06004 & Seminar hasil KKN: presentasi dan defense laporan serta solusi TI. & Modalitas: Blended Learning. Bentuk: Luring/Daring. Metode: presentasi dan Case Method. & 1 x 2 x 50' seminar hasil. & Mahasiswa mempresentasikan dan mempertahankan laporan KKN beserta solusi TI yang dihasilkan di hadapan tim penguji. & Bentuk penilaian: Seminar hasil (presentasi dan tanya jawab). Mengacu rubrik RTM. & Kualitas presentasi; ketepatan jawaban; kualitas solusi TI. & 25 \\
\end{longtblr}
\footnotetext[1]{\textit{Draf konsep untuk perbandingan. Nama mata kuliah dan bobot 2 SKS mengikuti Nota Dinas Wakil Direktur I No. 11/WADIR I/KM/2026. Kode mata kuliah, semester, dan integrasi ke crosswalk kurikulum akan difinalkan pada versi yang dipilih.}}
